%% kapa: exam=1 round=23 year=2012 period=1 form=A subject=economics subject_ko=경제원론 seq=7 number=47
%% source: Q-Net 감정평가사 기출문제 (cst00309) artlSeq=5052578 file=raw/1008/5052578/2012년 제23회 감정평가사 제1차시험 기출문제(1교시 A형).pdf
%% conversion: kordoc 4.24.51 -> tools/gen1.py (auto); review=proofread-sonnet; flags=none
\begin{question}{47}
기업 A의 노동과 자본의 투입량과 산출량 수준을 관찰한 결과 다음과 같은 표를 얻었다. 이 표에서 발견할 수 없는 현상은? (단, 생산에 투입되는 요소는 노동과 자본 뿐이다.)
\begin{center}
\renewcommand{\arraystretch}{1.1}
\begin{tabular}{|c|c|c|}\hline
노동투입 & 자본투입 & 총생산 \\ \hline
1 & 4 & 20 \\ \hline
2 & 2 & 20 \\ \hline
3 & 2 & 28 \\ \hline
4 & 1 & 20 \\ \hline
4 & 2 & 35 \\ \hline
4 & 3 & 38 \\ \hline
4 & 4 & 40 \\ \hline
\end{tabular}
\end{center}
\begin{choices}
  \item 규모의 경제
  \item 규모수익 불변
  \item 노동의 한계생산 체감
  \item 자본의 한계생산 체감
  \item 노동에 대한 자본의 한계대체율 체감
\end{choices}
\end{question}
\answer{1}
